\section{Finite-horizon scheduling and rolling implementation}
\label{sec:model}

The capacity tables from Section~\ref{sec:capacity} determine a time-expanded
MILP for fixed initial fleets and given order windows. The same formulation
is then used with revealed orders and inherited commitments in rolling execution.

\subsection{Decisions and objective}

For each applicable arc and direction, let $x_{ij}^{kd}$ count the human-driven
vehicles assigned to a local wave, $y_{ij}^d$ the automated vehicles,
$w_{ij}^d$ the human-driven vehicles assigned to a direct wave, and
$e_{ij}^d$ the independent empty transfers. All four are nonnegative integers. Binary selectors
$u_{ijv}^{kd}$ and $\widetilde u_{ijv}^d$ identify a positive local or
direct vehicle count $v$. Zero vehicles are represented by selecting no positive
size. For $\ell\in\mathcal L^d$, the nonnegative continuous variables
$g_{ij}^{k\ell}$, $\widehat g_{ij}^{\ell}$, and $h_{ij}^{\ell}$ are the
local, automated, and direct quantities of that order. The unserved quantity
is $0\leq z_\ell\leq D_\ell$.

Let $c_H,c_A,c_D,c_E\geq0$ be costs in currency per vehicle-minute.
When hourly rates are supplied to the implementation, they are divided by
$60$ once. With zero hub-transfer delay and no separate transfer fee, the
objective is
\begin{align}
 \min\quad &\sum_{\ell\in\mathcal L}p_\ell z_\ell
 +c_H\Delta t\sum_{k,d}\sum_{(i,j)\in\mathcal A_k}(j-i)x_{ij}^{kd}
 +c_A\Delta t\sum_d\sum_{(i,j)\in\mathcal A_A}(j-i)y_{ij}^d
 \nonumber\\
 &+c_D\Delta t\sum_d\sum_{(i,j)\in\mathcal A_D}(j-i)w_{ij}^d
 +c_E\Delta t\sum_d\sum_{(i,j)\in\mathcal A_E}(j-i)e_{ij}^d.
 \label{mod:objective}
\end{align}
The default empty-transfer cost is $c_E=c_D$, although it can be specified
separately. The unit of $p_\ell$ is currency per stop equivalent, so
both rejection and vehicle-time terms are monetary costs. The objective
prices vehicle operation and rejection; accumulation waiting, shipment
lead time, and sorting delay are outside this finite-horizon objective.
Vehicle cost is charged for the entire assigned task duration,
including any slack within that interval. Unserved-demand penalties allow
rejection, rather than imposing compulsory acceptance of every order.

\subsection{Wave and line-haul capacity}

The local constraints, for every $k,d,(i,j)\in\mathcal A_k$, are
\begin{equation}
 \begin{aligned}
 x_{ij}^{kd}&=\sum_{v=1}^N v u_{ijv}^{kd},\qquad
 \sum_{v=1}^N u_{ijv}^{kd}\leq1,\\
 \sum_{\ell\in\mathcal L^d}g_{ij}^{k\ell}
 &\leq\sum_{v=1}^N K^k_{ij}(v)u_{ijv}^{kd}.
 \end{aligned}
 \label{mod:local-links}
\end{equation}
For each $d,(i,j)\in\mathcal A_D$, impose
\begin{equation}
 \begin{aligned}
 w_{ij}^d&=\sum_{v=1}^N v\widetilde u_{ijv}^d,\qquad
 \sum_{v=1}^N\widetilde u_{ijv}^d\leq1,\\
 \sum_{\ell\in\mathcal L^d}h_{ij}^{\ell}
 &\leq\sum_{v=1}^N\widetilde K^d_{ij}(v)\widetilde u_{ijv}^d.
 \end{aligned}
 \label{mod:direct-links}
\end{equation}
Automated capacity is shared across all orders on an arc:
\begin{equation}
 \sum_{\ell\in\mathcal L^d}\widehat g_{ij}^{\ell}
 \leq M_A y_{ij}^d,
 \qquad d\in\mathcal D,\ (i,j)\in\mathcal A_A.
 \label{mod:auto-capacity}
\end{equation}
Both selector ranges use the total human-driven fleet $N$, not just its
initial city allocation $N_k$: a vehicle can migrate and subsequently join
a local wave. The inventory constraints below enforce the actual
availability in each city. Neither $e$ nor $y$ requires a positive customer
load. Nor does \eqref{mod:direct-links} force $w$ to carry cargo: a
$\tau$-period zero-cargo direct movement remains admissible under the above
convention. The separate variable $e$ makes empty relocation available
independently of customer candidates and permits a distinct cost or duration.

\subsection{Vehicle conservation in time and space}

Let $d_k$ be the direction leaving city $k$ and $\bar d_k$ its opposite.
For every $k$ and $t\in\mathcal T$, define the idle human-driven inventory
\begin{align}
 I^H_k(t)={}&N_k
 -\sum_{(i,j)\in\mathcal A_D:i\leq t}w_{ij}^{d_k}
 -\sum_{(i,j)\in\mathcal A_E:i\leq t}e_{ij}^{d_k}
 +\sum_{(i,j)\in\mathcal A_D:j\leq t}w_{ij}^{\bar d_k}
 +\sum_{(i,j)\in\mathcal A_E:j\leq t}e_{ij}^{\bar d_k}
 \nonumber\\
 &-\sum_{d\in\mathcal D}\sum_{(i,j)\in\mathcal A_k:i\leq t<j}x_{ij}^{kd}
 \ \geq0.
 \label{mod:human-inventory}
\end{align}
The corresponding automated inventory is
\begin{equation}
 I^A_k(t)=A_k-
 \sum_{(i,j)\in\mathcal A_A:i\leq t}y_{ij}^{d_k}
 +\sum_{(i,j)\in\mathcal A_A:j\leq t}y_{ij}^{\bar d_k}
 \geq0.
 \label{mod:auto-inventory}
\end{equation}
Local occupancy is half-open, $i\leq t<j$, so a vehicle completing at $j$
can start a new task at that same node. The same convention allows an
intercity arrival and a departure at the same node. Summing the two
automated inventory inequalities also gives the implied fleet bound
\begin{equation}
 \sum_d\sum_{(i,j)\in\mathcal A_A:i\leq t<j}y_{ij}^d\leq A.
 \label{mod:auto-busy}
\end{equation}
For human-driven vehicles, \eqref{mod:human-inventory} simultaneously
accounts for local occupancy, migration, and independent empty transfers;
a fixed per-city busy-vehicle bound alone would not do so.

\subsection{Order windows, precedence, and demand balance}

For $\ell\in\mathcal L^d$, impose the forbidden combinations
\begin{align}
 g_{ij}^{o(d)\ell}&=0 &&\text{if }i<s_\ell,
 &g_{ij}^{r(d)\ell}&=0 &&\text{if }j>e_\ell,
 \label{mod:local-windows}\\
 h_{ij}^{\ell}&=0 &&\text{if }i<s_\ell\text{ or }j>e_\ell.
 \label{mod:direct-windows}
\end{align}
Delivery is credited at the end of the destination tour or direct trip.
For each $t\in\mathcal T$, cumulative precedence requires
\begin{align}
 \sum_{(i,j)\in\mathcal A_A:i\leq t}\widehat g_{ij}^{\ell}
 &\leq\sum_{(i,j)\in\mathcal A_{o(d)}:j\leq t}g_{ij}^{o(d)\ell},
 \label{mod:pickup-before-auto}\\
 \sum_{(i,j)\in\mathcal A_{r(d)}:i\leq t}g_{ij}^{r(d)\ell}
 &\leq\sum_{(i,j)\in\mathcal A_A:j\leq t}\widehat g_{ij}^{\ell}.
 \label{mod:auto-before-delivery}
\end{align}
Goods can wait at either hub, but cannot depart before arrival. Equality
of event times permits immediate handover; no sorting queue is modeled in
this finite-horizon formulation. Finally, require
\begin{equation}
 \sum_{(i,j)\in\mathcal A_k}g_{ij}^{k\ell}
 +\sum_{(i,j)\in\mathcal A_D}h_{ij}^{\ell}+z_\ell=D_\ell,
 \qquad k\in\{o(d),r(d)\},\quad \ell\in\mathcal L^d.
 \label{mod:demand}
\end{equation}
At $t=T$, \eqref{mod:pickup-before-auto}--\eqref{mod:demand} imply that
the total transshipment quantity is identical at collection, automated
line-haul, and distribution. Consequently, an order cannot receive rejection
credit for goods that have already been collected and left unfinished.
The direct and transshipment quantities are separate portions of the same
demand and are not counted twice.

Equations \eqref{mod:objective}--\eqref{mod:demand}, the arc domains, binary
selectors, nonnegative integer vehicle counts, and nonnegative continuous
goods and rejection variables define a mixed-integer linear program once
the capacity tables have been computed. Transfer-only operation with city-dedicated human-driven fleets is the
restriction $w=e=h=0$; the direct-wave selectors are then zero as well.
This restriction closes both direct freight service and human-driven empty
intercity repositioning. Comparing it with the flexible model therefore
changes both freight routing and human-driven fleet mobility.
The core formulation places no prescribed lower or upper share on direct
service. Positive transfer delays, transfer fees, and direct-share bounds
are optional implementation extensions and are not used in the model
analyzed here.

\subsection{Service variants and the physical endpoint}

The default model leaves the final city allocation free. This is a distinct
choice from allowing unfinished tasks: every modeled arc still ends by
$T$. If an actual shutdown requires human-driven terminal stocks
$\bar N_1+\bar N_2=N$, it can instead impose
\begin{equation}
 I^H_k(T)=\bar N_k,\qquad k=1,2.
 \label{mod:optional-terminal}
\end{equation}
Such a constraint must account for both $w$ and $e$. A computational
rolling-window boundary is not an operating shutdown and does not justify
\eqref{mod:optional-terminal}. The rolling implementation described in
Section~\ref{sec:method}
uses the free terminal allocation; enforcing a genuine future terminal
target would additionally require a feasible completion policy for all
previous commitments.
